% group4_quantum_computing_synthesis.tex
% GeoVac Group 4 Synthesis: Quantum Computing / Structural Sparsity Arc
% Status: Project archive synthesis, June 2026
% Audience: quantum-computing / NISQ / VQE researchers
% Spine: encoding (14) -> resources (20) -> periodicity foundation (16) -> nuclear extension (23)

\documentclass[aps,pra,twocolumn,superscriptaddress,longbibliography]{revtex4-2}

\usepackage{amsmath,amssymb,amsthm,graphicx,xcolor,bm,braket,booktabs}
\usepackage{hyperref}

\newtheorem{theorem}{Theorem}
\newtheorem{proposition}{Proposition}
\newtheorem{corollary}{Corollary}
\newtheorem{observation}{Observation}


% ---------------------------------------------------------------
% \gvq{registry-key}{value} -- a load-bearing quantity, carrying a
% foreign key into debug/qa/numeric_registry.py alongside its literal
% value.  Renders as the value alone, so this file and the archived
% PDF stay self-contained; the key gives the QA gate (C20) referential
% integrity, and `--index <key>` names every locus citing the
% quantity when it moves.  See CLAUDE.md SS15.
\newcommand{\gvq}[2]{#2}

\begin{document}

\title{Structural Sparsity for Quantum Simulation:\\
A Synthesis of the GeoVac Quantum-Computing Arc}
\thanks{Group~4 synthesis in the Geometric Vacuum series}

\author{J.~Loutey}
\affiliation{Independent Researcher, Kent, Washington}

\date{June 28, 2026}

%% ========================================================================
%% ABSTRACT
%% ========================================================================

\begin{abstract}
The Geometric Vacuum (GeoVac) framework discretizes electronic
structure on natural geometries, producing sparse graph Hamiltonians
whose angular content is fixed by Gaunt selection rules.  This synthesis
assembles the framework's quantum-computing arc---the four Group~4
papers---into a single resource-oriented narrative.  The spine is:\ the
qubit encoding and its Pauli/$1$-norm scaling~\cite{loutey_paper14}; the
benchmarked resource library across three periodic-table
rows~\cite{loutey_paper20}; the $S_N$ representation-theoretic
foundation that classifies the atomic building
blocks~\cite{loutey_paper16}; and the nuclear shell-model extension that
delineates exactly which part of the framework
transfers~\cite{loutey_paper23}.  The unifying value proposition is
\emph{structural}, not accuracy-based:\ after the Jordan--Wigner
transformation the GeoVac lattice Hamiltonian yields $\sim Q^{3.8}$ Pauli
terms for atoms (corrected 2026-08-29; the retired pair-diagonal encoding
gave $O(Q^{3.15})$) and a Pauli count exactly linear in qubit count for
composed multi-center molecules at fixed basis (within-molecule basis
exponent $3.17$; corrected 2026-08-29 from the retired $O(Q^{2.5})$)
(versus $O(Q^{3.9\text{--}4.3})$ for Gaussian baselines), with a
sub-quadratic Pauli $1$-norm $\lambda = O(Q^{1.77})$ ($R^2=0.9997$;
re-measured 2026-08-30, retired: $O(Q^{1.69})$) that
controls Trotter depth and qubitization query cost.  Two disciplines run
through the arc.  \emph{First}, we separate what survived the
2026-08-29 selection-rule correction from what did not.  \emph{Surviving:}
the exact linearity $N_{\mathrm{Pauli}} = c\,Q$ across molecules at fixed
basis, the isostructural invariance, the $Z$-independence, and the
equal-qubit 1-norm advantage ($3.8\times$ at $Q=10$, $7.1\times$ at
$Q=28$) --- each re-measured under the exact rule and re-confirmed, not
inherited.  \emph{Not surviving:} the within-molecule and
atomic exponents (which moved $2.5 \to 3.17$ and $3.15 \to 3.8$), the QWC
exponent (which moved $3.36 \to 4.01$, above the term exponent, withdrawing
the measurement-group advantage entirely), and every absolute multiplier.  The earlier framing of this discipline --- ``scaling
exponents unchanged, only multipliers re-price'' --- was itself retracted
with the correction:\ the re-pricing is not a constant factor, it grows
along the basis axis.
\emph{Second}, every comparison is made at \emph{matched qubit count},
not matched accuracy:\ the composed basis carries classical
equilibrium-geometry errors of $5$--$26\%$, so the framework is
positioned for fixed-geometry NISQ/VQE simulation (Pauli count, QWC
measurement groups), not for fault-tolerant geometry optimisation.  The
nuclear extension is honest in the same way:\ its binding energies are
encoding-validation benchmarks far from experiment, and the Fock
projection rigidity theorem states precisely that the conformal $S^3$
structure does \emph{not} transfer to the harmonic-oscillator basis---only
the angular machinery and the qubit encoding do.
\end{abstract}

\maketitle

%% ========================================================================
\section{Introduction: a structural value proposition}
%% ========================================================================

The classical-solver investigation that occupies the GeoVac
quantum-chemistry arc (Group~2) reached a clear verdict:\ the
natural-geometry hierarchy characterises what discrete graph
Hamiltonians can and cannot achieve for ground-state potential-energy
surfaces, and on the accuracy axis it is not competitive with production
quantum chemistry.  The framework's computational value lies elsewhere.
The same construction that makes the classical accuracy bounded---a
single-center angular basis with Gaunt-enforced selection rules and a
pseudopotential-partitioned composed geometry---makes the resulting
\emph{qubit} Hamiltonians structurally sparse.  This synthesis collects
the four papers that develop and benchmark that observation.

The reader should hold two framings in mind throughout.  First, the
advantage is in \emph{resource scaling and term counts}, measured at
\emph{matched qubit count}, not in energy accuracy:\ at $n_{\max}=2$ the
composed basis carries $5$--$26\%$ classical equilibrium-geometry errors
(LiH/BeH$_2$/H$_2$O), so the natural use is single-point simulation at a
fixed nuclear geometry, where the sparsity is free and the geometry
error does not enter.  Second, the framework targets the NISQ/VQE
regime---where Pauli term count and the number of qubitwise-commuting
(QWC) measurement groups dominate cost---rather than the fault-tolerant
QPE regime, although the sub-quadratic Pauli $1$-norm is also the
relevant quantity for qubitization.

The four papers form a dependency spine.
Paper~14~\cite{loutey_paper14} establishes the encoding and its scaling
laws.  Paper~20~\cite{loutey_paper20} benchmarks a $37$-system library
(35 composed molecules + He + H$_2$) across three periodic-table rows and exports them to the
standard quantum-software ecosystems.  Paper~16~\cite{loutey_paper16}
supplies the $S_N$ representation-theoretic classification of the atomic
building blocks that the composed builder assembles.
Paper~23~\cite{loutey_paper23} extends the angular machinery to the
nuclear shell model and proves, via the Fock rigidity theorem, the
boundary of what transfers.

%% ========================================================================
\section{The angular-sparsity mechanism}
%% ========================================================================

The structural origin of every resource result in this arc is a single
fact about the two-electron integral (ERI) tensor.  In a single-center
basis indexed by $(n,\ell,m)$, the angular part of each ERI is a product
of Gaunt coefficients, and a Gaunt coefficient vanishes unless its three
angular momenta satisfy the $3j$ triangle inequality and the magnetic
selection rule.  As the number of spatial orbitals $M$ grows, the
fraction of nonzero ERIs therefore \emph{decays}---empirically as
$\approx M^{-0.49}$ (corrected 2026-08-29; the retired rule's
$M^{-0.85}$ overstated the decay)---whereas a Gaussian ERI tensor remains near-dense at
$O(1)$ occupancy.  This is the potential-independent angular-sparsity
theorem of Paper~22 (Group~3):\ the ERI density depends only on
$\ell_{\max}$, not on the radial potential $V(r)$.  The sparsity
survives the Jordan--Wigner transformation because it is carried by the
angular labels, which JW does not mix.

\paragraph{The selection rule, and a correction to it.}
The exact Coulomb angular selection rule is the global-$M_L$ rule,
$m_a+m_b=m_c+m_d$, and every number in this synthesis is now computed
under it.

Until 2026-08-29 the production composed and atomic builders realised
instead a \emph{pair-diagonal} restriction $m_a=m_c \wedge m_b=m_d$,
described in earlier versions of this document as a deliberate sparsifying
approximation.  An independent re-derivation traced it to a \emph{sign
error}:\ the Gaunt coefficient set the Wigner-3j magnetic argument to
$q = m_c-m_a$ where the bottom row requires $q = m_a-m_c$, so the
coefficient vanished identically unless $m_a=m_c$, silently deleting every
$m$-changing multipole (88 of the 126 nonzero $c^k$ at $l\le2$) --- among
them the genuinely nonzero ``m-swap'' integrals
(e.g.\ $\braket{p_{+1}p_{-1}|p_{-1}p_{+1}}$, angular factor
$\sqrt{6}/5 \approx 0.49$).  The error was arbitrated against direct
quadrature of the angular factor from its definition.

Two consequences, both of which retract earlier statements here:
\begin{itemize}
\item The re-pricing is \emph{not} a constant factor.  It is
$2.51\times$ uniformly across main-group molecules at $n_{\max}=2$, but
grows along the basis axis ($1.00\times/2.51\times/5.40\times$ at
$n_{\max}=1/2/3$), so the within-molecule exponent moves $2.54 \to 3.17$.
The earlier ``constant factor, scaling unchanged'' claim was a
single-point artifact.
\item Two defects were corrected together and they behave oppositely,
which is worth separating because the distinction was previously
collapsed.  The wrong-sign Gaunt argument ($q = m_c - m_a$) deleted every
$m$-changing multipole and so changed the ERI \emph{support}:\ 65
$\rightarrow$ \gvq{block_sp_n2_eri}{107} nonzero entries per block, which is why every count,
density and exponent moved.  The Condon--Shortley factor-order error, by
contrast, leaves the ERI support exactly invariant and flips signs only
--- measured on the relativistic block as identical support (2{,}676
entries) and identical $\sum|{\rm ERI}|$ with $40.4\%$ of signs reversed.
It is not resource-neutral even so:\ the sign flips change which
Jordan--Wigner terms cancel, moving LiH$_{\rm rel}$ from $1{,}413$ to
$\gvq{lih_rel_n2_pauli}{1{,}501}$ Pauli terms.  The distinction is between the ERI level, where
only the sign error moved the support, and the resource level, where both
defects moved counts.  ``Values,
not entries'' is true of the second and false of the first; an earlier
version of this synthesis generalised it to the whole correction and
thereby claimed the counts were invariant when they were not.  What
survives is the \emph{form} $N_{\mathrm{Pauli}} = c\,Q$, with its
coefficient re-measured ($11.10 \rightarrow \gvq{composed_coeff}{27.90}$).
\end{itemize}
Every resource number reported here is the exact global-$M_L$ value;
retired pair-diagonal figures appear only where explicitly marked as
retired history.  There is no dual-rule convention to adopt:\ the
pair-diagonal rule was a wrong-sign Gaunt argument, not an alternative to
the exact one.

%% ========================================================================
\section{The encoding and its scaling laws (Paper~14)}
%% ========================================================================

Paper~14 transforms the GeoVac lattice Hamiltonian to qubits via
Jordan--Wigner and characterises the resulting Pauli decomposition by
several operationally distinct metrics.

\paragraph{Term count.}
For single-center atoms the Pauli term count scales as $\sim Q^{3.8}$ in
the qubit number $Q$ (fit $3.78$; corrected 2026-08-29 from the retired
$O(Q^{3.15})$), which sits at the lower edge of --- not below --- the
Gaussian molecular Pauli-scaling exponents in the $Q^{3.9}$--$Q^{4.3}$
range; the
specific electronic-structure values $Q^{4.25}$ (LiH) and $Q^{3.92}$
(H$_2$O), and the range they bracket, are GeoVac's own log-log fit of
the published Gaussian Pauli counts of Trenev et al.~\cite{trenev2025}
(Table~5, Appendix~B; Jordan--Wigner with 2-qubit reduction).
For composed multi-center molecules the pseudopotential-partitioned
architecture produces block-diagonal ERIs and a \emph{universal} within-molecule exponent of $3.17$ (identical to all
reported digits --- corrected 2026-08-29 from the retired $O(Q^{2.5})$,
whose apparent $0.02$ spread was an artifact of the sign error), across the
benchmarked molecules---independent of molecular topology.  Across the
library the composed count is, remarkably, \emph{exactly linear} in the
qubit count,
\begin{equation}
N_{\mathrm{Pauli}} = c\,Q,\qquad
c = 27.90~(\text{main group}),\quad 30.03~(d\text{-block}),
\label{eq:linearity}
\end{equation}
under the exact global-$M_L$ rule at $n_{\max}=2$.  (The retired
pair-diagonal rule gave $11.10$ and $9.23$; note that the ordering
\emph{reverses} --- the $d$-block, which appeared \emph{sparser} under the
retired rule, is in fact \emph{denser}, because higher angular momentum
carries more $m$-swap channels for the sign error to drop.  The
pair-diagonal $d$-block sparsity was an artifact, not a selection-rule
property.)

\paragraph{Measurement and Trotter cost.}
Two further metrics govern practical cost, and the 2026-08-30
re-measurement separates them sharply.  The number of QWC measurement
groups---the dominant factor in the measurement overhead of variational
algorithms---scales as $O(Q^{4.01})$, which is \emph{steeper} than the
Pauli term exponent ($3.77$) and lies inside the Gaussian
$O(Q^{4\text{--}5})$ band.  No measurement-group scaling advantage is
therefore claimed; under the retired pair-diagonal rule this exponent read
$O(Q^{3.36})$, below the term count, and that reading is withdrawn.  The
Pauli $1$-norm $\lambda = \sum_i |c_i|$, which bounds both Trotter step
counts and qubitization query complexity, scales as $O(Q^{1.77})$ with
$R^2 = 0.9997$ (retired: $O(Q^{1.69})$).  This sub-quadratic $1$-norm
scaling survives the correction and remains the key
fault-tolerant-relevant result:\ Trotter circuit depth grows far more
slowly with system size than the raw term count alone would suggest.
The \emph{molecular} Gaussian baselines (LiH, H$_2$O) are the published
Pauli counts of Trenev et al.~\cite{trenev2025} (Table 5, Appendix B),
cross-checked against analytic references (Szabo and Ostlund for H$_2$;
analytical Slater-type orbitals for He).  The 6-31G and cc-pVDZ entries of
Paper~14's Table~I use synthetically generated integrals and are retained
for scaling-fit purposes only.

\paragraph{Tier discipline.}
The scaling exponents are \emph{measured} quantities (log-log fits with
reported $R^2$), not derived; the angular selection rules underlying the
ERI-density decay are \emph{exact}; and the composed sparsity advantage
\emph{magnitude} is a measured quantity under a named ERI rule.  The
construction itself is zero-parameter:\ the Hamiltonian follows from
nuclear charges and geometry alone.

%% ========================================================================
\section{The benchmarked resource library (Paper~20)}
%% ========================================================================

Paper~20 turns the encoding into a usable artifact:\ the
\texttt{hamiltonian()} registry of $37$ systems ($35$ composed molecular
qubit Hamiltonians spanning $Z=1$--$56$ (H through Ba),
plus the He and H$_2$ reference atoms), including multi-center systems
(CO, N$_2$, F$_2$), mixed frozen-core molecules (NaCl), and
the first transition series as hydrides (ScH through ZnH), exported to
Qiskit, OpenFermion, and PennyLane.

\paragraph{Two architectures.}
The library exposes two composed-molecular architectures with
complementary strengths.  The \emph{composed} (pseudopotential, PK)
architecture is the sparser of the two:\ LiH requires $\gvq{lih_composed_pauli}{838}$ Pauli terms
at $30$ qubits, with $4.3\times$ fewer QWC groups ($64$ vs.\ $273$) than the
Gaussian STO-3G reference at $12$ qubits (raw baseline).  The small-system
market test is near parity on the raw-vs-raw comparison --- $838$ versus
STO-3G's $\gvq{lih_sto3g_raw}{907}$ ($1.08\times$ fewer) --- and \emph{unfavourable} against the
qubit-reduced count ($276$, i.e.\ GeoVac is $3.0\times$ larger).
(Corrected 2026-08-29; the retired pair-diagonal figures were $334$ Pauli
and $13\times$ QWC.)  The sparsity advantage spans $0.28\times$ (H$_2$O against minimal-basis
STO-3G, where GeoVac is \emph{denser}) through $55$--$76\times$ against
cc-pVDZ to $317\times$ at the most favourable equal-qubit point; the retired
pair-diagonal rule supported a blanket ``one to three orders of magnitude''
that survives at neither end.  What advantage there is rests on the scaling
gap and the multi-center block structure, \emph{not} on the small-molecule
minimal-basis count---which is near parity.  The \emph{balanced coupled} (PK-free)
architecture trades sparsity for chemistry:\ it binds LiH with an FCI
minimum at a computed $R_{\mathrm{eq}}=3.227$~bohr (7.0\% above the
experimental 3.015~bohr) at $n_{\max}=2$ ($\gvq{lih_balanced_pauli}{2{,}726}$ Pauli terms, $30$
qubits; corrected 2026-08-29 from the retired $878$);
a single-point evaluation at the experimental geometry reaches a $0.20\%$
energy error at $n_{\max}=3$ when solved exactly (number-projected FCI).

\paragraph{Row-conditional chemistry scope.}
The chemistry-accuracy scope at $n_{\max}=2$ is honestly
\emph{row-conditional}:\ first-row hydrides (LiH, BeH$_2$, H$_2$O, \dots)
bind with documented PES topology, while second-row hydrides (NaH and
below) exhibit monotone overattraction under every tested
configuration---frozen-core balanced, multi-zeta valence, screened
cross-center, explicit-core Hartree--Fock---with no interior PES minimum
at the qubit-Hamiltonian level.  This scope boundary is documented
empirically, not papered over:\ the framework's value proposition is
structural sparsity with a proven state-space Gromov--Hausdorff
truncation rate (Paper~38) --- a \emph{metric-space} rate, explicitly
\emph{not} a direct energy-error bound (the operator-to-energy lift
overshoots by 3--4 orders of magnitude) --- not energy-accuracy parity with
CCSD(T) on second-row systems.

\paragraph{Frozen cores.}
Second- and third-row molecules use analytical frozen cores ([Ne], [Ar],
[Ar]$3d^{10}$); the frozen-core energy enters only through the identity
operator and does not affect the quantum measurement cost.
Isostructural molecules with frozen cores produce identical Pauli counts
across rows---the periodic-table structure is visible directly in the
resource counts.

%% ========================================================================
\section{The representation-theoretic foundation (Paper~16)}
%% ========================================================================

The composed builder assembles molecules from atomic building blocks,
and Paper~16 supplies the group-theoretic classification of those
blocks.  After separation of the hyperradius, the $N$-electron angular
dynamics live on $S^{3N-1}$ and the angular kinetic energy is the
$\mathrm{SO}(3N)$ Casimir operator with eigenvalue
$\mu_{\mathrm{free}} = \nu(\nu + 3N - 2)/2$.  Fermionic antisymmetry
restricts the spatial wavefunction to the $S_N$ irreducible
representation conjugate to the spin Young diagram, whose first-column
length $\lambda_1$ fixes the minimum angular quantum number
$\nu = N - \lambda_1$.  For every atom with ground-state spin $S<N/2$ the
spatial irrep has $\lambda_1=2$, giving the universal result
\begin{equation}
\nu = N-2,\qquad \mu_{\mathrm{free}} = 2(N-2)(N-1).
\end{equation}

This produces a five-type classification of atoms---A (trivial), B
(noble gas), C (alkali), D (alkaline earth), E (open shell)---and recasts
the periodic law as the repetition of the $S_N$ irrep sequence
$\mathrm{C}\to\mathrm{D}\to\mathrm{E}\to\mathrm{B}$ within each period.
The classification is implemented in \texttt{atomic\_classifier.py},
which refines the open-shell type~E into a $p$-block dispatch (E) and a
separate $d$-block dispatch (F) for the composed builder, and supplies
the universal angular quantum number $\nu = N-2$.

\paragraph{Honest framing.}
The classification \emph{maps onto, rather than predicts}, the known
periodic table:\ its value is a unified group-theoretic language for a
pattern usually described empirically, not a derivation of chemistry
from first principles.  Paper~16 also establishes a structural fact
relevant to the relativistic extension:\ the topology $(\nu,
\mu_{\mathrm{free}})$ is smooth through $Z = 1/\alpha \approx 137$, so
the Dirac instability is a \emph{metric} singularity on $S^{3N-1}$, not a
topological one.

%% ========================================================================
\section{The nuclear extension and its boundary (Paper~23)}
%% ========================================================================

Paper~23 applies the angular-sparsity machinery to a system the
conformal $S^3$ projection cannot reach---the nuclear shell model---and
in doing so delineates exactly which part of the framework is universal.

\paragraph{What the angular machinery delivers.}
Built on the harmonic-oscillator single-particle basis with a
Mayer--Jensen spin--orbit term, the construction reproduces the HO shell
closures $2, 8, 20, 40, 70, 112$ directly from graph state counting, and
recovers the physical magic numbers $2, 8, 20, 28, 50, 82, 126$ with a
spin--orbit strength $v_{ls}/\hbar\omega \approx 0.17$ and a Nilsson
$\ell(\ell+1)$ correction $d_{\ell\ell}/\hbar\omega \approx 0.02$.  The
qubit encodings inherit the angular sparsity:\ the deuteron ($1p+1n$)
requires $16$ qubits and $688$ non-identity Pauli terms using the
Minnesota nucleon--nucleon potential and
Moshinsky--Talmi brackets; helium-4 ($2p+2n$) requires $16$ qubits and
$828$ Pauli terms despite a Hilbert space $12.25\times$ larger---only
$1.20\times$ more terms, a direct consequence of the angular selection
rules carrying through to multi-particle systems.  A composed
nuclear--electronic deuterium encoding ($26$ qubits) validates the
multi-block architecture and quantifies the $\sim 10^{13}$ coefficient
ratio between nuclear, electronic, and hyperfine scales.

\paragraph{The rigidity boundary.}
\begin{theorem}[Fock projection rigidity]
The $S^3$ conformal projection that makes electronic Hamiltonians
equivalent to a graph Laplacian is unique to the Coulomb potential and
does not extend to the harmonic-oscillator or Woods--Saxon systems.
\end{theorem}
For nuclear Hamiltonians the GeoVac contribution is therefore the
angular machinery and qubit encoding only, \emph{not} the conformal
projection.  This is the structural dual of the foundations-side HO
rigidity theorem (Paper~24).

\paragraph{Honest framing.}
All nuclear physics inputs (Minnesota parameters, spin--orbit strength,
$\hbar\omega$) follow the standard shell-model literature---they are not
GeoVac-derived---and the quantitative binding energies at
$N_{\mathrm{shells}}=2$ are far from experiment.  They are to be read as
\emph{encoding-validation benchmarks}, demonstrating that the two-species
sparse qubit architecture works, not as a nuclear-structure calculation.

%% ========================================================================
\section{Honest scope and positioning}
%% ========================================================================

\paragraph{What is robust.}
What survives the correction is \emph{structure}, not numbers.  The
exact linearity $N_{\mathrm{Pauli}} = c\,Q$ across molecules at fixed
basis \eqref{eq:linearity}, the isostructural invariance, the
$Z$-independence of the ERI count, and the block-diagonal
selection-rule sparsity all hold under the exact rule, and rest on the
potential-independent angular-sparsity theorem (Paper~22).  Indeed the
linearity is \emph{cleaner} after the correction:\ the two-point
within-molecule exponent is now identical across all six molecules and
forced rather than fitted.

Every \emph{exponent}, by contrast, moved:\ atomic $3.15 \to 3.8$,
within-molecule $2.5 \to 3.17$, $1$-norm $1.69 \to 1.774$, QWC
$3.36 \to \gvq{exp_qwc_3pt}{4.013}$ --- and the QWC exponent moved far enough to overturn its
own conclusion, since $\gvq{exp_qwc_3pt}{4.013}$ exceeds the Pauli exponent and lies inside
the Gaussian $O(Q^{4\text{--}5})$ band.  An earlier version of this
paragraph listed those exponents under ``robust'' while stating two
sentences later that they had moved.  The qubit
library is a concrete, exported, reproducible artifact.

\paragraph{What is conditional.}
The \emph{absolute} sparsity multipliers and small-system market-test
lines were re-priced on 2026-08-29 when the pair-diagonal ERI restriction
was found to be a sign error; the re-pricing is basis-dependent, not a
constant factor, and it leaves composed LiH at near parity with raw
STO-3G ($838$ vs.\ $\gvq{lih_sto3g_raw}{907}$) and $3.0\times$ larger than the qubit-reduced
count.  All Pauli comparisons are at matched qubit count, not
matched accuracy:\ the composed basis carries $5$--$26\%$ classical
$R_{\mathrm{eq}}$ errors, the chemistry-accuracy scope is row-conditional
(first-row binds; second-row overattracts), and the nuclear binding
energies are encoding-validation benchmarks far from experiment.

\paragraph{What was tried and did not work.}
The arc includes several honest negatives that bound the design space.
Discrete $\mathbb{Z}_2$ tapering beyond the Hopf $m_\ell \to -m_\ell$
symmetry does not transfer to relativistic chemistry:\ neither
$\kappa$-parity, direct-basis $m_j$-parity, nor rotated-basis
$m_j$-parity commutes with the Dirac--Coulomb Hamiltonian, because the
$jj$-coupled Coulomb operator couples sign-counted sectors under
sum-conservation (not parity) rules.  Adding off-diagonal cross-block
$h_1$ coupling to gain binding produces $16\times$ over-binding on LiH;
the framework's binding wall --- the structural obstruction to
chemistry-accuracy binding documented in the Group~2 record --- is
structurally inseparable from its
encoding advantages.  A non-abelian gauge upgrade of the composed Dirac
does not reduce the Pauli count below the $\mathbb{Z}_2$-tapered
identity, because tapering requires a Hamiltonian \emph{symmetry} and the
gauge freedom is a basis change, not a symmetry.  The spectral action is
the wrong functional for chemistry observables at finite cutoff.  These
negatives are recorded so they are not re-derived.

\paragraph{Positioning.}
The natural target is variational, NISQ-era simulation---where the Pauli
term count and the number of QWC measurement groups dominate cost---at
fixed molecular geometry.  The sub-quadratic Pauli $1$-norm makes the
encodings relevant to qubitization as well, but the framework is not
positioned as a fault-tolerant-QPE or geometry-optimisation tool, and it
is not a replacement for production quantum chemistry.  It is a
structurally sparse, zero-parameter, ecosystem-exported library of qubit
Hamiltonians with a proven state-space Gromov--Hausdorff truncation rate
(a metric-space rate, not an energy-error bound)---a research
instrument for the resource-estimation and algorithm-development side of
quantum simulation.

%% ========================================================================
%% BIBLIOGRAPHY
%% ========================================================================
\begin{thebibliography}{99}

\bibitem{loutey_paper7}
J.~Loutey,
``The Dimensionless Vacuum:\ Recovering the Schr\"odinger Equation from Scale-Invariant Graph Topology,''
GeoVac Paper~7 (2026).

\bibitem{loutey_paper14}
J.~Loutey,
``Structurally Sparse Qubit Hamiltonians from Spectral Graph Theory,''
GeoVac Paper~14 (2026).

\bibitem{loutey_paper16}
J.~Loutey,
``The Periodic Table Recast as $S_N$ Representation Theory on $S^{3N-1}$,''
GeoVac Paper~16 (2026).

\bibitem{loutey_paper20}
J.~Loutey,
``Structurally Sparse Molecular Hamiltonians for Near-Term Quantum Simulation,''
GeoVac Paper~20 (2026).

\bibitem{loutey_paper22}
J.~Loutey,
``Potential-Independent Angular Sparsity for Qubit Hamiltonians in Spherical Fermion Systems,''
GeoVac Paper~22 (2026).

\bibitem{loutey_paper23}
J.~Loutey,
``Nuclear Shell Model Hamiltonians on the Hyperspherical Lattice,''
GeoVac Paper~23 (2026).

\bibitem{trenev2025}
D.~Trenev, P.~J.~Ollitrault, S.~M.~Harwood, T.~P.~Gujarati,
S.~Raman, A.~Mezzacapo, and S.~Mostame,
``Refining resource estimation for the quantum computation of
vibrational molecular spectra through Trotter error analysis,''
Quantum \textbf{9}, 1630 (2025); arXiv:2311.03719.
[Cited for the electronic-structure Gaussian-basis Pauli counts in
Table~5 (Appendix~B; Jordan--Wigner with 2-qubit reduction), the
Gaussian baseline; the $Q^{3.9}$--$Q^{4.3}$ scaling exponents (and the
specific LiH/H$_2$O exponents) are GeoVac's own log-log fit of those
published counts, per Paper~14.]

\end{thebibliography}

\end{document}
